% GENERATED FILE. Edit canonical lessons, then run npm run generate:books.
% canonical-source-hash: fnv1a64:5a9b7ad76802227a
% canonical-lessons: GU-C530-avdhi, GU-C530-takid, GU-C530-mahetal, GU-C530-thelvu, GU-C530-akhre

\chapter{Deadlines}
\label{ch:gu-deadlines}
\hlchaptermodality{530}

\begin{chapteropening}
\textbf{By the end of this chapter:} \emph{I can talk about deadlines, urgent work and extra time.}

Complete the last lesson of chapter 530: I can talk about deadlines, urgent work and extra time.
\end{chapteropening}

\section[\texorpdfstring{\gu{અવધિ} (avdhi)}{અવધિ (avdhi)}]{\gu{અવધિ} (avdhi) — a time limit}

\label{lesson:GU-C530-avdhi}



\begin{quote}
Before the new one: say the Gujarati for half-finished, then the Gujarati for final.
\end{quote}

\subsection*{You'll want to know: \gu{અવધિ}}
\textbf{\gu{અવધિ}} — \emph{avdhi} — \textquotedblleft{}a time limit\textquotedblright{}.

\begin{quote}
\gu{જો} \gu{તમે} \gu{આ} \gu{અવધિમાં} \gu{અરજી} \gu{નહીં} \gu{કરો} \gu{તો} \gu{તમે} \gu{આ} \gu{તક} \gu{ગુમાવશો}.
\end{quote}

\begin{quote}
\emph{jo tame ā avdhimā̃ arjī nahī̃ karo to tame ā tak gumāvśo.}
\end{quote}

\begin{quote}
If you don't apply within this time limit, you will lose this chance.
\end{quote}

A feminine noun.

\begin{cousinweb}
\textbf{\gu{અવધિ}} (\emph{avdhi}) is a tatsama, taken from Sanskrit \emph{avadhi}, \textquotedblleft{}a term, a limit\textquotedblright{} (Monier-Williams).
\end{cousinweb}

\begin{grammarlens}[title={What you've built}]
One more word for saying what happens if.
\end{grammarlens}

\subsection*{Your turn}
\begin{itemize}
\raggedright
  \item \emph{Say it:} \emph{avdhi}
  \item \emph{Say it:} \emph{avdhi}, once more
  \item \emph{Say it:} \emph{jo tame ā avdhimā̃ arjī nahī̃ karo to tame ā tak gumāvśo.} — If you don't apply within this time limit, you will lose this chance.
  \item \emph{Recall:} say the Gujarati for half-finished, then the Gujarati for final, then say \emph{avdhi} again
\end{itemize}

\begin{tcolorbox}[breakable,colback=teal!4,colframe=teal!35!black,title={Before you move on}]
What does \gu{અવધિ} mean? (\textquotedblleft{}A time limit\textquotedblright{}.) Say it once more.
\end{tcolorbox}
\section[\texorpdfstring{\gu{તાકીદ} (tākīd)}{તાકીદ (tākīd)}]{\gu{તાકીદ} (tākīd) — urgency}

\label{lesson:GU-C530-takid}



\begin{quote}
Before the new one: say the Gujarati for final, then the Gujarati for a time limit.
\end{quote}

\subsection*{You'll want to know: \gu{તાકીદ}}
\textbf{\gu{તાકીદ}} — \emph{tākīd} — \textquotedblleft{}urgency\textquotedblright{}.

\begin{quote}
\gu{જો} \gu{કામમાં} \gu{તાકીદ} \gu{હશે} \gu{તો} \gu{હું} \gu{રવિવારે} \gu{પણ} \gu{આવીશ}.
\end{quote}

\begin{quote}
\emph{jo kāmmā̃ tākīd haśe to hũ ravivāre paṇ āvīś.}
\end{quote}

\begin{quote}
If the work is urgent, I'll come in on Sunday too.
\end{quote}

A feminine noun.

\begin{cousinweb}
\textbf{\gu{તાકીદ}} (\emph{tākīd}) is usually traced to Arabic \emph{taʾkīd}, \textquotedblleft{}emphasis, insistence\textquotedblright{}, through Persian; in Gujarati it is the pressing need to get something done.
\end{cousinweb}

\begin{grammarlens}[title={What you've built}]
One more word for saying what happens if.
\end{grammarlens}

\subsection*{Your turn}
\begin{itemize}
\raggedright
  \item \emph{Say it:} \emph{tākīd}
  \item \emph{Say it:} \emph{tākīd}, once more
  \item \emph{Say it:} \emph{jo kāmmā̃ tākīd haśe to hũ ravivāre paṇ āvīś.} — If the work is urgent, I'll come in on Sunday too.
  \item \emph{Recall:} say the Gujarati for final, then the Gujarati for a time limit, then say \emph{tākīd} again
\end{itemize}

\begin{tcolorbox}[breakable,colback=teal!4,colframe=teal!35!black,title={Before you move on}]
What does \gu{તાકીદ} mean? (\textquotedblleft{}Urgency\textquotedblright{}.) Say it once more.
\end{tcolorbox}
\section[\texorpdfstring{\gu{મહેતલ} (mahetal)}{મહેતલ (mahetal)}]{\gu{મહેતલ} (mahetal) — extra time}

\label{lesson:GU-C530-mahetal}



\begin{quote}
Before the new one: say the Gujarati for a time limit, then the Gujarati for urgency.
\end{quote}

\subsection*{You'll want to know: \gu{મહેતલ}}
\textbf{\gu{મહેતલ}} — \emph{mahetal} — \textquotedblleft{}extra time\textquotedblright{}.

\begin{quote}
\gu{જો} \gu{સાહેબ} \gu{મને} \gu{એક} \gu{અઠવાડિયાની} \gu{મહેતલ} \gu{આપશે} \gu{તો} \gu{હું} \gu{કામ} \gu{પતાવી} \gu{દઈશ}.
\end{quote}

\begin{quote}
\emph{jo sāheb mane ek aṭhvāḍiyānī mahetal āpśe to hũ kām patāvī daīś.}
\end{quote}

\begin{quote}
If the boss gives me a week's extra time, I'll get the work finished.
\end{quote}

A feminine noun.

\begin{cousinweb}
\textbf{\gu{મહેતલ}} (\emph{mahetal}) is usually traced to Arabic \emph{muhla}, \textquotedblleft{}a respite, a delay granted\textquotedblright{}, through Persian.
\end{cousinweb}

\begin{grammarlens}[title={What you've built}]
One more word for saying what happens if.
\end{grammarlens}

\subsection*{Your turn}
\begin{itemize}
\raggedright
  \item \emph{Say it:} \emph{mahetal}
  \item \emph{Say it:} \emph{mahetal}, once more
  \item \emph{Say it:} \emph{jo sāheb mane ek aṭhvāḍiyānī mahetal āpśe to hũ kām patāvī daīś.} — If the boss gives me a week's extra time, I'll get the work finished.
  \item \emph{Recall:} say the Gujarati for a time limit, then the Gujarati for urgency, then say \emph{mahetal} again
\end{itemize}

\begin{tcolorbox}[breakable,colback=teal!4,colframe=teal!35!black,title={Before you move on}]
What does \gu{મહેતલ} mean? (\textquotedblleft{}Extra time\textquotedblright{}.) Say it once more.
\end{tcolorbox}
\section[\texorpdfstring{\gu{ઠેલવું} (ṭhelvũ)}{ઠેલવું (ṭhelvũ)}]{\gu{ઠેલવું} (ṭhelvũ) — to push back (a date)}

\label{lesson:GU-C530-thelvu}



\begin{quote}
Before the new one: say the Gujarati for urgency, then the Gujarati for extra time.
\end{quote}

\subsection*{You'll want to know: \gu{ઠેલવું}}
\textbf{\gu{ઠેલવું}} — \emph{ṭhelvũ} — \textquotedblleft{}to push back (a date)\textquotedblright{}.

\begin{quote}
\gu{જો} \gu{વરસાદ} \gu{પડશે} \gu{તો} \gu{અમે} \gu{પ્રવાસ} \gu{આવતા} \gu{અઠવાડિયા} \gu{પર} \gu{ઠેલીશું}.
\end{quote}

\begin{quote}
\emph{jo varsād paḍśe to ame pravās āvtā aṭhvāḍiyā par ṭhelīśũ.}
\end{quote}

\begin{quote}
If it rains, we will push the trip back to next week.
\end{quote}

Past, with \gu{તેણે} (he or she …): \gu{તેણે} \gu{પ્રવાસ} \gu{ઠેલ્યો}, \emph{teṇe pravās ṭhelyo}.

\begin{cousinweb}
\textbf{\gu{ઠેલવું}} (\emph{ṭhelvũ}) comes from \emph{*ṭhēḍḍ}, a form CDIAL reconstructs, and CDIAL lists the Gujarati verb as \textquotedblleft{}to push\textquotedblright{} (CDIAL 5512); a plan or a date is pushed back.
\end{cousinweb}

\begin{grammarlens}[title={What you've built}]
One more word for saying what happens if.
\end{grammarlens}

\subsection*{Your turn}
\begin{itemize}
\raggedright
  \item \emph{Say it:} \emph{ṭhelvũ}
  \item \emph{Say it:} \emph{ṭhelvũ}, once more
  \item \emph{Say it:} \emph{jo varsād paḍśe to ame pravās āvtā aṭhvāḍiyā par ṭhelīśũ.} — If it rains, we will push the trip back to next week.
  \item \emph{Recall:} say the Gujarati for urgency, then the Gujarati for extra time, then say \emph{ṭhelvũ} again
\end{itemize}

\begin{tcolorbox}[breakable,colback=teal!4,colframe=teal!35!black,title={Before you move on}]
What does \gu{ઠેલવું} mean? (\textquotedblleft{}To push back (a date)\textquotedblright{}.) Say it once more.
\end{tcolorbox}
\section[\texorpdfstring{\gu{આખરે} (ākhre)}{આખરે (ākhre)}]{\gu{આખરે} (ākhre) — at long last}

\label{lesson:GU-C530-akhre}



\begin{quote}
Before the new one: say the Gujarati for extra time, then the Gujarati for to push back.
\end{quote}

\subsection*{You'll want to know: \gu{આખરે}}
\textbf{\gu{આખરે}} — \emph{ākhre} — \textquotedblleft{}at long last\textquotedblright{}.

\begin{quote}
\gu{જો} \gu{આપણે} \gu{મહેનત} \gu{કરીશું} \gu{તો} \gu{આખરે} \gu{આપણે} \gu{જીતીશું}.
\end{quote}

\begin{quote}
\emph{jo āpaṇe mahenat karīśũ to ākhre āpaṇe jītīśũ.}
\end{quote}

\begin{quote}
If we work hard, in the end we will win.
\end{quote}

\begin{cousinweb}
\textbf{\gu{આખરે}} (\emph{ākhre}) is usually traced to Arabic \emph{ākhir}, \textquotedblleft{}last\textquotedblright{}, through Persian, with the Gujarati ending \emph{-e} of \textbf{\gu{ગઈકાલે}}: at the last.
\end{cousinweb}

\begin{grammarlens}[title={What you've built}]
One more word for saying what happens if.
\end{grammarlens}

\subsection*{Your turn}
\begin{itemize}
\raggedright
  \item \emph{Say it:} \emph{ākhre}
  \item \emph{Say it:} \emph{ākhre}, once more
  \item \emph{Say it:} \emph{jo āpaṇe mahenat karīśũ to ākhre āpaṇe jītīśũ.} — If we work hard, in the end we will win.
  \item \emph{Recall:} say the Gujarati for extra time, then the Gujarati for to push back, then say \emph{ākhre} again
\end{itemize}

\begin{tcolorbox}[breakable,colback=teal!4,colframe=teal!35!black,title={Before you move on}]
What does \gu{આખરે} mean? (\textquotedblleft{}At long last\textquotedblright{}.) Say it once more.
\end{tcolorbox}
